
\documentclass{article}
\usepackage{colortbl}
\usepackage{makecell}
\usepackage{multirow}
\usepackage{supertabular}

\begin{document}

\newcounter{utterance}

\centering \large Interaction Transcript for game `cladder', experiment `full\_v1.5\_default', episode 1296 with Qwen3.6{-}27B.
\vspace{24pt}

{ \footnotesize  \setcounter{utterance}{1}
\setlength{\tabcolsep}{0pt}
\begin{supertabular}{c@{$\;$}|p{.15\linewidth}@{}p{.15\linewidth}p{.15\linewidth}p{.15\linewidth}p{.15\linewidth}p{.15\linewidth}}
    \# & \multicolumn{2}{c}{Player} && \multicolumn{2}{c}{Game Master} \\
    \hline

    \theutterance \stepcounter{utterance}  
    & & & \multicolumn{4}{p{0.6\linewidth}}{
        \cellcolor[rgb]{0.9,0.9,0.9}{
            \makecell[{{p{\linewidth}}}]{
                \texttt{\tiny{[P1$\langle$GM]}}
                \texttt{You are an expert in causal inference. The following question is not a typical commonsense query, but rather a meticulously designed question created by a professor specializing in causal inference, intended to assess the students' mastery of the course content.} \\
\texttt{Imagine a self{-}contained, hypothetical world with only the following conditions, and without any unmentioned factors or causal relationships: Gender has a direct effect on smoking and lung cancer. Smoking has a direct effect on tar deposit. Tar deposit has a direct effect on lung cancer. Gender is unobserved. The overall probability of smoking is 12\%. For nonsmokers, the probability of lung cancer is 58\%. For smokers, the probability of lung cancer is 61\%. Is lung cancer more likely than absence of lung cancer overall?} \\
\texttt{Start your answer with 'Yes' or 'No', followed by additional reasoning or evidence to support your explanation.} \\
            }
        }
    }
     \\ \\

    \theutterance \stepcounter{utterance}  
    & \multicolumn{4}{p{0.6\linewidth}}{
        \cellcolor[rgb]{0.561,0.737,0.561}{
            \makecell[{{p{\linewidth}}}]{
                \texttt{\tiny{[P1$\rangle$GM]}}
                \texttt{No} \\
\\ 
\texttt{To determine whether lung cancer is more likely than its absence overall, we must calculate the total probability of lung cancer, denoted as \$P(L)\$, using the law of total probability. We are given the marginal probability of smoking, \$P(S) = 0.12\$, and the conditional probabilities of lung cancer given smoking status: \$P(L|S) = 0.61\$ and \$P(L|\textbackslash{}neg S) = 0.58\$.} \\
\\ 
\texttt{The calculation is as follows:} \\
\texttt{\$\$P(L) = P(L|S)P(S) + P(L|\textbackslash{}neg S)P(\textbackslash{}neg S)\$\$} \\
\\ 
\texttt{First, determine the probability of not smoking:} \\
\texttt{\$\$P(\textbackslash{}neg S) = 1 {-} P(S) = 1 {-} 0.12 = 0.88\$\$} \\
\\ 
\texttt{Next, substitute the values into the total probability formula:} \\
\texttt{\$\$P(L) = (0.61 \textbackslash{}times 0.12) + (0.58 \textbackslash{}times 0.88)\$\$} \\
\texttt{\$\$P(L) = 0.0732 + 0.5104\$\$} \\
\texttt{\$\$P(L) = 0.5836\$\$} \\
\\ 
\texttt{The overall probability of lung cancer is \$0.5836\$ (or \$58.36\textbackslash{}\%\$). Consequently, the probability of the absence of lung cancer is \$1 {-} 0.} \\
            }
        }
    }
    & & \\ \\

    \theutterance \stepcounter{utterance}  
    & & & \multicolumn{2}{p{0.3\linewidth}}{
        \cellcolor[rgb]{0.9,0.9,0.9}{
            \makecell[{{p{\linewidth}}}]{
                \texttt{\tiny{[GM$|$GM]}}
                \texttt{No} \\
            }
        }
    }
    & & \\ \\

    \theutterance \stepcounter{utterance}  
    & & & \multicolumn{2}{p{0.3\linewidth}}{
        \cellcolor[rgb]{0.9,0.9,0.9}{
            \makecell[{{p{\linewidth}}}]{
                \texttt{\tiny{[GM$|$GM]}}
                \texttt{yes} \\
            }
        }
    }
    & & \\ \\

    \theutterance \stepcounter{utterance}  
    & & & \multicolumn{2}{p{0.3\linewidth}}{
        \cellcolor[rgb]{0.9,0.9,0.9}{
            \makecell[{{p{\linewidth}}}]{
                \texttt{\tiny{[GM$|$GM]}}
                \texttt{game\_result = LOSE} \\
            }
        }
    }
    & & \\ \\

\end{supertabular}
}

\end{document}
